\section{Modus Ponens Reverse}

\bigskip

\begin{definition}
\label{bij_Fin}
\lean{bij_Fin}
\uses{def:Z2n}
\leanok
For $m\in\mathbb{N}$, let $g : \Fin (2^m) \to \Zz^m$ be a bijective function. $g$ exists, {\B applying} the fact that $\Fin (2^m)$ and $\Zz^m$ {\B simply} have {\B equal cardinals}.
\end{definition}

\verb|Mathlib hint:| \href{https://leanprover-community.github.io/mathlib4_docs/Mathlib/Data/Fintype/EquivFin.html#Fintype.equivOfCardEq}{Fintype.equivOfCardEq}.

\bigskip

\begin{lemma}
\label{sum_basis}
\lean{sum_basis}
\uses{def:e}
\leanok
If $i \in \Fin n$ and $f : \Fin n \to \Zz^m$, then $\displaystyle\sum_{j \in \Fin n} e_{ij} \boldsymbol{\cdot} f(j) = f(i).$
\end{lemma}

\begin{proof}
  \leanok
{\B Unfolding} the definition of \eqref{e}, we can {\B rewrite}  $\sum_{j \in \Fin n} e_{ij} \boldsymbol{\cdot} f(j)$ as $f(i)$, because the {\B single} term that contributes in the {\B sum} is the one with $j = i$ and $e_{ii} = 1$.
\end{proof}

\verb|Mathlib hint:| \href{https://leanprover-community.github.io/mathlib4_docs/Mathlib/Algebra/Module/BigOperators.html#Fintype.sum_single_smul}{Fintype.sum\_single\_smul}.

\bigskip

\begin{theorem}
\label{puzzle_mpr}
\lean{puzzle_mpr}
\uses{def:SF}
\leanok
If $n = 2^m$ for some $m$, then the puzzle has a strategy for $n$.
\end{theorem}

\begin{proof}
\uses{bij_Fin, sum_basis}
\leanok
First, we {\B rewrite} $n$ as $2^m$ in the definition of \eqref{SF}. {\B Let} $g$ be given by {\B Definition} \ref{bij_Fin} and {\B use} for $f$ in \eqref{SF} the function
$$b\mapsto g^{-1}\left(\sum_{j\in\Fin 2^m} b_j \boldsymbol{\cdot} g(j)\right).$$

{\B Introduce} $b\in\Zz^n$ and $k\in\Fin n$ in order to verify the conditions of \eqref{SF}. We want to show
$$\exists\,i\in\Fin n, f(b + {\rm e}_i) = g^{-1}\left(\sum_{j\in\Fin 2^m} (b_j + e_{ij}) \boldsymbol{\cdot} g(j)\right) = k.$$

We {\B simplify} the above relation using the {\B distributivity of the scalar multiplication over addition} and over {\B sum} to get
$$\sum_{j\in\Fin 2^m} (b_j + e_{ij}) \boldsymbol{\cdot} g(j) = \sum_{j\in\Fin 2^m} b_j \boldsymbol{\cdot} g(j) + \sum_{j\in\Fin 2^m} e_{ij} \boldsymbol{\cdot} g(j).$$

{\B Let}'s {\B use} $i := g^{-1}\left(g(k) - \sum_{j\in\Fin 2^m} b_j \boldsymbol{\cdot} g(j)\right)$.
By {\B Lemma \ref{sum_basis}}, we can {\B rewrite} the desired relation with $\displaystyle\sum_{j\in\Fin 2^m} e_{ij} \boldsymbol{\cdot} g(j) = g(i)$. By the expression of $i$, we can {\B simplify} to get
\begin{align*}
& g^{-1}\left(\sum_{j\in\Fin 2^m} b_j \boldsymbol{\cdot} g(j) + g(i)\right) \\
&= g^{-1}\left(\sum_{j\in\Fin 2^m} b_j \boldsymbol{\cdot} g(j) + g(k) - \sum_{j\in\Fin 2^m} b_j \boldsymbol{\cdot} g(j)\right) = g^{-1}(g(k)) = k.
\end{align*}
\end{proof}

\verb|Mathlib hints:| \href{https://leanprover-community.github.io/mathlib4_docs/Mathlib/Logic/Equiv/Defs.html#Equiv.symm}{g.symm}, \href{https://leanprover-community.github.io/mathlib4_docs/Mathlib/Algebra/Module/Defs.html#add_smul}{add\_smul}, \href{https://leanprover-community.github.io/mathlib4_docs/Mathlib/Algebra/BigOperators/Group/Finset/Basic.html#Finset.sum_add_distrib}{Finset.sum\_add\_distrib}.
